%\begin{bibunit}
%\nocite{*}
%\putbib[Biblio]
%\end{bibunit}



%%%%%%%%%%%%%%%%%%%%%%%%%%%%%%%%%%%%%%%%%%%%%%%%%%%%%%%%%%%%%%
%%%%%%%%%%%%% MA-XXXX Matemática para el Cálculo %%%%%%%%%%%%%
%%%%%%%%%%%%%%%%%%%%%%%%%%%%%%%%%%%%%%%%%%%%%%%%%%%%%%%%%%%%%%
\newpage 
\subsection*{MA-01XX Matemática para el Cálculo}
\addcontentsline{toc}{subsection}{MA-01XX Matemática para el Cálculo}
\begin{refsection}

\begin{table}[H]
\centering{
\begin{tabular}{|ll|}
\hline
%\multicolumn{2}{|c|}{\textbf{Nivel de virtualidad del curso:} Virtual}\\ \hline
\textbf{Año:} I  & \textbf{Requisitos:} Ninguno \\ 
\textbf{Ciclo:} I  & \textbf{Correquisitos:} Ninguno \\ 
\textbf{Tipo de curso:} Teórico & \textbf{Horas presenciales por semana:}   \\ 
\textbf{Créditos:} 4 & \textbf{Horas de trabajo independiente por semana:}  \\ \hline
\end{tabular}
}
\end{table}

\subsubsection*{Descripción del curso}

Este curso es una transición entre el nivel de matemática que ha adquirido el estudiantado durante sus años de enseñanza media y un nivel de matemática superior, requerido en cursos posteriores. Entendemos matemática como conceptos y entes abstractos, sus propiedades y las relaciones entre estos. Dada la relevancia de esta comprensión, el curso se ubica en primer ciclo de carrera, acompañado por el curso MA-0100 Introducción al Cálculo, conformando el primer acercamiento a un nivel de matemática más avanzado.

El curso ha sido diseñado para iniciar este acercamiento, retomando aprendizajes previos con el fin de reforzar y estructurar el desarrollo de habilidades operacionales. Estos conocimientos brindarán herramientas para el óptimo desempeño en cursos posteriores. Este curso se complementará con MA-2XXX Matemática Introductoria II.

\subsubsection*{Objetivos}

\begin{enumerate}
\item Reproducir procedimientos en los contenidos matemáticos del curso que luego permitan al estudiantado interpretar, descubrir, inferir y hasta predecir argumentaciones y conexiones dentro de la matemática.
\item Estimular (fomentar) en el estudiantado capacidades de análisis e interpretación que permitan el desarrollo de habilidades operacionales y el primer estímulo de pensamiento abstracto, ambos conceptos piedras angulares presentes en los mecanismos de razonamiento posteriores.
\end{enumerate}

\subsubsection*{Contenidos}

\begin{enumerate}
\item \textbf{Números racionales y reales:} Cuerpos ordenados, ecuaciones e inecuaciones polinomiales  y paramétricas. (2 semanas)

\item \textbf{Geometría analítica y trigonometría:} Geometría analítica básica: coordenadas cartesianas (en el plano?), distancia euclideana, lugares geométricos, ecuaciones cartesianas de círculos y cónicas, transformaciones (traslaciones). Trigonometría básica: Las funciones trigonométricas y sus gráficas, identidades trigonométricas, coordenadas polares, rotaciones, ....

\item \textbf{Teoría elemental de números:} Divisibilidad básica en $\mathbb{Z}$, el algoritmo de la división, máximo común divisor, mínimo común múltiplo, coprimalidad, la identidad de Bézout, números primos, factorización única en $\mathbb{Z}$, ternas pitagóricas, congruencias, ecuaciones lineales en congruencias, el Teorema Pequeño de Fermat y el Teorema de Euler. (4 semanas).

\item \textbf{Números complejos:} Definiciones y propiedades básicas, parte real e imaginaria, conjugación, valor absoluto, representación polar de números complejos, la fórmula de De Moivre, potencias, raíces $n$-ésimas, el Teorema Fundamental del Álgebra, factorización de polinomios en $\mathbb{R}$ y en $\mathbb{C}$. (2 semanas).

\item \textbf{Inducción matemática:} principio de inducción matemática, sumatorias, productorias, definiciones recursivas, ejemplos y prácticas. (2 semanas).
\end{enumerate}

\subsubsection*{Metodología}

\subsubsection*{Bibliografía}
El libro de Axler %\cite{axler2011algebra}

\printbibliography[heading=subbibliography,section=\the\value{refsection}]
\end{refsection}
